\section*{अध्याय \ref{ch:b1:kinetic-theory} --- अणुगति सिद्धांत और आदर्श गैस}

\begin{solution}{exo:b1:kinetic-theory:1}
$N = PV/k_BT = 1.013\times10^5 \times 10^{-3}/(1.38\times10^{-23} \times 273) =
\num{2.7e22}$; $n^{*-1/3} = (3.7\times10^{-26})^{1/3} = \qty{3.3}{nm}$, यानी
आणविक आकार का दस गुना।
\end{solution}

\begin{solution}{exo:b1:kinetic-theory:2}
$u = \sqrt{3RT/M}$: H$_2$ \qty{1930}{m/s}, N$_2$ \qty{517}{m/s}, O$_2$
\qty{484}{m/s}, CO$_2$ \qty{412}{m/s}; \qty{77}{K} पर N$_2$: \qty{262}{m/s}।
\end{solution}

\begin{solution}{exo:b1:kinetic-theory:3}
निरपेक्ष $P = \qty{3.5}{bar}$: $n = PV/RT = 3.5\times10^5 \times 0.030/(8.314
\times 293) = \qty{4.3}{mol}$, यानी $\qty{125}{g}$। \qty{323}{K} पर, $P \propto T$:
\qty{3.86}{bar} (प्रमापी \qty{2.85}{bar})।
\end{solution}

\begin{solution}{exo:b1:kinetic-theory:4}
$n^* = P/k_BT = \qty{2.4e25}{m^{-3}}$; $\lambda = 1/(\sqrt2\pi \times 1.37\times
10^{-19} \times 2.4\times10^{25}) = \qty{6.8e-8}{m}$; आवृत्ति $u/\lambda \approx
500/7\times10^{-8} = \qty{7e9}{s^{-1}}$। $\lambda \propto 1/P$: \qty{1}{m} तब,
जब $P \approx 10^5 \times 7\times10^{-8} = \qty{7e-3}{Pa}$ हो।
\end{solution}

\begin{solution}{exo:b1:kinetic-theory:5}
$\rho = PM/RT$: \qty{293}{K} पर \qty{1.20}{kg/m^3}, \qty{373}{K} पर
\qty{0.946}{kg/m^3}। उत्थापन $(1.20 - 0.95) \times 2800 \times 9.81 = \qty{7.0}{kN}$:
यानी लगभग \qty{720}{kg}।
\end{solution}

\begin{solution}{exo:b1:kinetic-theory:6}
$P_{\mathrm{O}_2} = 0.21P$: \qty{0.21}{atm}; \qty{5000}{m} पर \qty{0.11}{atm};
और एवरेस्ट पर \qty{0.07}{atm} --- यानी समुद्र-तल का एक-तिहाई। ऑक्सीजन रक्त
में अपने आंशिक \omterm{def:b1:kinetic-theory:pressure}{दाब} के समानुपात में घुलती है; एक-तिहाई पर फेफड़े
हीमोग्लोबिन को संतृप्त नहीं कर पाते।
\end{solution}

\begin{solution}{exo:b1:kinetic-theory:7}
\omterm{def:b1:kinetic-theory:model}{आदर्श गैस}: $P = RT/V = 2494/10^{-3} = \qty{24.9}{bar}$। वान डेर वाल्स:
$RT/(V - b) - a/V^2 = 2494/9.57\times10^{-4} - 0.364/10^{-6} = 26.1 - 3.6 =
\qty{22.4}{bar}$: आकर्षण ($-3.6$ \omterm{def:b1:kinetic-theory:pressure}{बार}) वर्जित आयतन ($+1.2$ \omterm{def:b1:kinetic-theory:pressure}{बार}) पर भारी पड़ता
है; और वास्तविक CO$_2$ आदर्श मान से $10\%$ नीचे रहती है।
\end{solution}

\begin{solution}{exo:b1:kinetic-theory:8}
$n = PV/RT = 10^5 \times 50/(8.314 \times 300) = \qty{2000}{mol}$; $U = \tfrac52
nRT = \qty{1.25e7}{J}$; $\Delta U = \tfrac52 nR \times 5 = \qty{208}{kJ}$;
\qty{2}{kW} पर \qty{104}{s}।
\end{solution}

\begin{solution}{exo:b1:kinetic-theory:9}
$\Delta L = \alpha L\Delta T = 1.2\times10^{-5} \times 30 \times 50 = \qty{18}{mm}$।
छल्ला: $\Delta d/d = \alpha\Delta T = 0.05/49.95 = 10^{-3}$: $\Delta T = 10^{-3}/
1.9\times10^{-5} = \qty{53}{K}$ (छेद भी अपने चारों ओर की धातु की तरह ही फैलता है)।
\end{solution}

\begin{solution}{exo:b1:kinetic-theory:10}
$\tfrac32 k_BT = \qty{6.2e-21}{J} = \qty{0.039}{eV}$। आर्गन $\tfrac32 R =
\qty{12.5}{J/(mol.K)}$; नाइट्रोजन $\tfrac52 R = \qty{20.8}{J/(mol.K)}$:
द्विपरमाणुक अणु घूर्णन में भी ऊर्जा संचित करता है (दो अतिरिक्त
स्वातंत्र्य-कोटियाँ), चाहे उसका द्रव्यमान कुछ भी हो।
\end{solution}

\begin{solution}{exo:b1:kinetic-theory:11}
छह-दिशा प्रतिरूप: पाठ की तरह ही, $P = \tfrac13 n^*mu^2$। समदैशिकता:
$\langle v_x^2\rangle = \langle v_y^2\rangle = \langle v_z^2\rangle$ और इनका
योग $\langle v^2\rangle$ है, अतः हर एक $\tfrac13\langle v^2\rangle$। किसी
वितरण के साथ अभिवाह-तर्क $P = n^*m\langle v_x^2\rangle$ देता है (दीवार से
$v_x > 0$ वाले अणु टकराते हैं, हर एक $2mv_x$ संवेग देता है, और प्रति
चाल-वर्ग अभिवाह $n^*v_x/2$, जिस पर औसत लिया जाता है), अर्थात् $\tfrac13 n^*m\langle v^2\rangle$ --- वही।
\end{solution}

\begin{solution}{exo:b1:kinetic-theory:12}
H$_2$: \qty{1.9}{km/s}, O$_2$: \qty{0.48}{km/s}; पृथ्वी \qty{11.2}{km/s},
चंद्रमा \qty{2.4}{km/s}। पृथ्वी पर O$_2$ के लिए $v_e/u = 23$: उससे आगे
वितरण की पूँछ नितांत नगण्य है; H$_2$ के लिए $v_e/u \approx 6$:
हर वर्ष एक अत्यंत छोटा पर लगातार अंश निकल भागता है, और अरबों वर्षों में
हाइड्रोजन चली जाती है। चंद्रमा पर $v_e/u$ घटकर $5$ (O$_2$) और
$1.2$ (H$_2$) रह जाता है: सब कुछ रिसकर निकल गया।
\end{solution}

\begin{solution}{pb:b1:kinetic-theory:1}
\textbf{1.} $V = \qty{200}{m^3}$: $n = PV/RT = 1.013\times10^5 \times 200/(8.314
\times 293) = \qty{8320}{mol}$; $N = nN_A = \num{5.0e27}$।

\textbf{2.} $m = nM = \qty{241}{kg}$: यानी तीन आदमी।

\textbf{3.} $n^* = N/V = \qty{2.5e25}{m^{-3}}$; दूरी $n^{*-1/3} = \qty{3.4}{nm}$,
यानी नौ आणविक व्यास।

\textbf{4.} N$_2$: $\sqrt{3 \times 8.314 \times 293/0.028} = \qty{511}{m/s}$;
O$_2$: \qty{478}{m/s}। प्रति अणु वही $\tfrac32 k_BT$: यानी $\tfrac12 mu^2$ नियत,
अतः $u \propto 1/\sqrt m$।

\textbf{5.} $\tfrac32 k_BT = \qty{6.1e-21}{J}$; कुल $N \times 6.1\times10^{-21}
= \qty{3.0e7}{J}$; \qty{100}{km/h} पर \qty{1000}{kg} की कार: \qty{3.9e5}{J} ---
यानी कमरे की आणविक गति अस्सी कारों जितनी है।

\textbf{6.} $\lambda = 1/(\sqrt2\pi \times 1.37\times10^{-19} \times 2.5\times10^{25})
= \qty{66}{nm}$; समय $\lambda/u = 66\times10^{-9}/500 = \qty{1.3e-10}{s}$।

\textbf{7.} $S$ पर प्रति सेकंड टक्करें: $\tfrac16 n^*Su = \tfrac16 \times 2.5\times
10^{25} \times 10^{-4} \times 500 = \qty{2.1e23}{s^{-1}}$।

\textbf{8.} हर एक $2mu = 2 \times 4.8\times10^{-26} \times 500 = \qty{4.8e-23}{kg.m/s}$ देती है
($m = M/N_A = \qty{4.8e-26}{kg}$): बल $2.1\times10^{23} \times 4.8\times10^{-23}
= \qty{10}{N}$ प्रति \qty{1}{cm^2}: यानी \qty{1.0e5}{Pa}। प्रतिरूप वायुमंडलीय
\omterm{def:b1:kinetic-theory:pressure}{दाब} को फिर से दे देता है।

\textbf{9.} $F = PS = 1.013\times10^5 \times 20 = \qty{2.0e6}{N}$ --- यानी दो सौ
टन; वही \omterm{def:b1:kinetic-theory:pressure}{दाब} दूसरी ओर से भी लगता है।

\textbf{10.} $P' = P \times 303/293 = \qty{1.048e5}{Pa}$; परिणामी $\Delta P\,S = 3.5
\times10^3 \times 20 = \qty{70}{kN}$ बाहर की ओर: यानी दीवार पर सात टन।

\textbf{11.} \omterm{def:b1:kinetic-theory:pressure}{दाब} बराबर होने तक हवा बाहर बहती है: $n \propto 1/T$, स्थिर
$P$, $V$ पर: $8320 \times (1 - 293/303) = \qty{275}{mol}$, यानी \qty{8}{kg}
बाहर चली जाती है।

\textbf{12.} $\rho = PM/RT$: \qty{293}{K} पर \qty{1.205}{kg/m^3}; \qty{373}{K} पर
\qty{0.946}{kg/m^3}।

\textbf{13.} गरम हवा $0.946 \times 2800 \times 9.81 = \qty{26.0}{kN}$; हटाई गई
ठंडी हवा $1.205 \times 2800 \times 9.81 = \qty{33.1}{kN}$; उत्थापन \qty{7.1}{kN}
(\qty{725}{kg})।

\textbf{14.} $725 - 300 = \qty{425}{kg}$: यानी पाँच यात्री।

\textbf{15.} उत्थापन $\propto \rho_c - \rho_h = \rho_c(1 - 293/T)$: और
$1 - 293/373 = 0.215$ को दुगुना करने के लिए $T = 293/0.57 = \qty{514}{K}$ (\qty{240}{\degreeCelsius}) चाहिए,
जो नायलॉन की सहन-सीमा (लगभग \qty{120}{\degreeCelsius}) से कहीं ऊपर है।

\textbf{16.} हीलियम $\rho = 1.205 \times 4/29 = \qty{0.166}{kg/m^3}$: उत्थापन
$(1.205 - 0.166) \times 2800 \times 9.81 = \qty{28.5}{kN}$ (\qty{2.9}{t}), यानी चार
गुना अधिक --- पर हीलियम महँगी है और हर अवतरण पर खो जाती है, जबकि गरम हवा
मुफ़्त है और उसका उत्थापन ज्वालक से घटाया-बढ़ाया जा सकता है।

\textbf{17.} $0.21 \times 1.013\times10^5 = \qty{2.1e4}{Pa}$।

\textbf{18.} $n^*_{\mathrm{O}_2} = P_{\mathrm{O}_2}/k_BT$: कमरा $2.1\times10^4/
(1.38\times10^{-23} \times 293) = \qty{5.3e24}{m^{-3}}$; एवरेस्ट $0.21 \times 0.33
\times 1.013\times10^5/(1.38\times10^{-23} \times 250) = \qty{2.0e24}{m^{-3}}$:
यानी कमरे का $38\%$।

\textbf{19.} $H = RT/Mg = 8.314 \times 260/(0.029 \times 9.81) = \qty{7.6}{km}$;
$P_0/2$ पर $H\ln 2 = \qty{5.3}{km}$।

\textbf{20.} $c_s = \sqrt{1.4 \times 8.314 \times 293/0.029} = \qty{343}{m/s}$,
जबकि मिश्रण के लिए $u = \qty{503}{m/s}$: $c_s/u = \sqrt{\gamma/3} = 0.68$।
ध्वनि अणुओं के द्वारा ही ले जाया गया विक्षोभ है, अतः वह उन्हीं की कोटि की
चाल से चलती है --- कुछ कम, क्योंकि केवल संचरण की दिशा वाला घटक ही मायने
रखता है।

\textbf{21.} $c_s \propto \sqrt T$: $343\sqrt{303/293} = \qty{349}{m/s}$।
हीलियम: $\sqrt{(5/3) \times 8.314 \times 293/0.004} = \qty{1010}{m/s}$: स्वर-नली
की अनुनाद-आवृत्तियाँ $c_s$ के साथ बढ़ती हैं, अतः हर स्वरक तीन गुना ऊँचा हो
जाता है।

\textbf{22.} अधिक ठंडी: $343\sqrt{250/293} = \qty{317}{m/s}$, यानी $8\%$ धीमी;
\omterm{def:b1:kinetic-theory:pressure}{दाब} इसमें आता ही नहीं।

\textbf{23.} वही $N$ (वही $P$, $V$, $T$); द्रव्यमान $\times 40/29 = \qty{333}{kg}$;
चालें $\times\sqrt{29/40}$: $u = \qty{428}{m/s}$; और $U = \tfrac32 nRT$, $\tfrac52$ की
जगह: यानी $60\%$।

\textbf{24.} $U = \tfrac52 \times 8320 \times 8.314 \times 293 = \qty{5.1e7}{J}$;
$\Delta U = \tfrac52 nR \times 10 = \qty{1.7e6}{J}$।

\textbf{25.} एक ही औसत, $\langle v^2\rangle$, \omterm{def:b1:kinetic-theory:pressure}{दाब} को तय करता है
($\tfrac13 n^*m\langle v^2\rangle$), \omterm{def:b1:kinetic-theory:temperature}{ताप} को परिभाषित करता है
($\tfrac12 m\langle v^2\rangle = \tfrac32 k_BT$), अवस्था समीकरण और \omterm{def:b1:kinetic-theory:U}{आंतरिक
ऊर्जा} दे देता है, और ध्वनि की चाल का पैमाना भी नियत कर देता है।
\end{solution}
